\section{1. The Big Picture: What Is This
Repository?}\label{the-big-picture-what-is-this-repository}

This repository is a \textbf{private, structured research laboratory}
for systematically replicating 17 academic papers on imbalanced time
series forecasting and regression --- the core empirical work of a
Master's thesis at CIn--UFPE.

\subsection{1.1 What ``Replication'' Means
Here}\label{what-replication-means-here}

You are \textbf{not} re-implementing papers from scratch. The workflow
is:

\begin{enumerate}
\def\labelenumi{\arabic{enumi}.}
\tightlist
\item
  Obtain the \textbf{original authors' source code} and place it inside
  the repository.
\item
  \textbf{Run that code} in a clean, isolated virtual environment.
\item
  \textbf{Compare your outputs} against the results reported in the
  paper.
\item
  \textbf{Document everything}: what matched, what differed, and why.
\end{enumerate}

\begin{quote}
\small Your role is that of a \textbf{scientific auditor} --- not a
reimplementer.
\end{quote}

This distinction matters because it defines what work is automated and
what is always manual.

\subsection{1.2 The Automation Boundary}\label{the-automation-boundary}

The repository's design explicitly separates bookkeeping (automated)
from scientific judgment (always manual).

{\def\LTcaptype{none} % do not increment counter
\small
\begin{longtable}[]{@{}
  >{\raggedright\arraybackslash}p{(\linewidth - 2\tabcolsep) * \real{0.5000}}
  >{\raggedright\arraybackslash}p{(\linewidth - 2\tabcolsep) * \real{0.5000}}@{}}
\toprule\noalign{}
\begin{minipage}[b]{\linewidth}\raggedright
Automated by scripts
\end{minipage} & \begin{minipage}[b]{\linewidth}\raggedright
Always done manually by you
\end{minipage} \\
\midrule\noalign{}
\endhead
\bottomrule\noalign{}
\endlastfoot
Creating the experiment folder structure & Cloning or copying the
original source code \\
Replacing placeholder text in template files & Downloading datasets \\
Rebuilding the progress dashboard in \texttt{README.md}
& Running the experiments \\
Syncing paper metadata from the Excel spreadsheet & Interpreting results
and filling the comparison table \\
Generating this guide as a PDF & Marking an experiment as Done \\
\end{longtable}
}

The most important rule: \textbf{you control all data and code
insertion}. Scripts only handle administrative bookkeeping.

\subsection{1.3 Repository Goals}\label{repository-goals}

By the end of the thesis, this repository should contain:

\begin{itemize}
\tightlist
\item
  \textbf{17 fully replicated experiments}, each with a complete
  scientific report.
\item
  \textbf{Documented deviations} from original paper results, with
  explanations.
\item
  A \textbf{paper registry} of all 43 papers in the reading list, with
  replication status.
\item
  \textbf{Shared utilities} for metrics and visualization that are
  reusable across experiments.
\end{itemize}

\subsection{1.4 Reading This Guide}\label{reading-this-guide}

This guide is organized in the order you will need the information:

{\def\LTcaptype{none} % do not increment counter
\small
\begin{longtable}[]{@{}
  >{\raggedright\arraybackslash}p{(\linewidth - 2\tabcolsep) * \real{0.5000}}
  >{\raggedright\arraybackslash}p{(\linewidth - 2\tabcolsep) * \real{0.5000}}@{}}
\toprule\noalign{}
\begin{minipage}[b]{\linewidth}\raggedright
Chapter
\end{minipage} & \begin{minipage}[b]{\linewidth}\raggedright
When you need it
\end{minipage} \\
\midrule\noalign{}
\endhead
\bottomrule\noalign{}
\endlastfoot
2 --- Directory Structure & First session: orienting yourself \\
3 --- Experiment Lifecycle & Every experiment: knowing what stage you
are in \\
4 --- Core Workflow & Every session: the exact step-by-step protocol \\
5 --- Automation Scripts & When running or understanding the scripts \\
6 --- File Reference & When filling in any experiment file \\
7 --- Environments & When setting up a new Python environment \\
8 --- Git Workflow & After each work session \\
9 --- Adding Experiments & When adding experiment \#018 or beyond \\
10 --- Rules & Quick reference cheatsheet \\
11 --- Troubleshooting & When something breaks \\
\end{longtable}
}
